\documentclass[11pt]{article}
\usepackage{geometry}
\geometry{
	a4paper,
	margin=2cm,
	left=3cm}
\usepackage[utf8]{inputenc}
%\usepackage[german]{babel}
\usepackage[T1]{fontenc}
\usepackage{graphicx}

\setlength{\parindent}{0em} %setzt Paragraphen-Einzug auf 0

\title{Firmenporträt Kleinprogramm}

\begin{document}
\includegraphics[scale=0.9]{banner.pdf}

\section*{Kleinprogramm GmbH}

Im Jahr 2014 gründete Michael Leopold die Kleinprogramm GmbH. Kurz darauf nahm er seinen ehemaligen Kommilitonen Emanuel Günther als Entwickler mit in das Unternehmen auf. Seitdem wurden bereits mehrere Software-Projekte für unterschiedliche mittelständische Unternehmen realisiert.\\

Unsere Geschäftsräume befinden sich im Zentrum Dresdens und sind damit gut erreichbar.\\
Als Kleinprogramm GmbH ist es uns wichtig, Programme mit hohen Qualitäts- und Sicherheitsstandards zu entwickeln. Wir pflegen stets eine enge Zusammenarbeit mit unseren Kunden, um eine optimale Software-Lösung gewährleisten zu können. Wir bieten Leistungen in unterschiedlichen Bereichen von Smartphone-Apps bis Web-Anwendung. \\ \\

\textbf{Von Kleinunternehmer bis Großkonzern - wir finden für Sie die passende Lösung!} \\ \\

Um uns bei der weiteren Ausbildung von Studenten und Studentinnen zu engagieren,  arbeiten wir seit 2016 in enger Kooperation mit der Hochschule für Technik und Wirtschaft (HTW) Dresden zusammen. Als Praxispartner bieten wir Studenten und Studentinnen die Möglichkeit ein Praktikum bei uns zu absolvieren oder an einem interessanten Projekt mitzuarbeiten. Durch diese Zusammenarbeit wird weiterhin eine Mitwirkung in aktuellen regionalen Forschungs\-projekten ermöglicht. \\

Im Jahr 2019 wurden zwei weitere Mitarbeiter in das Unternehmen aufgenommen. Robert Siedler ist für Controlling und Finanzen zuständig, während Patrick Creutzburg als Werksstudent in unterschiedlichen Bereichen tätig ist.\\

\vspace*{\fill}
\includegraphics[scale=0.9]{footer.pdf}
\thispagestyle{empty}
\end{document}